\section{Methods}
\label{sec:methods}

\subsection{Part 1: generating the missing reason labels}

\paragraph{Data and features.}
 The Part~1 files, derived from
ClinicalTrials.gov~\cite{clinicaltrials} with some processing performed by course staff, hold 1,561 trials that stopped early: 1,171 training trials of
which only 157 carry visible labels (the remaining 1,014 have all five reasons
masked as $-100$), 78 fully labeled validation trials, and 312 unlabeled test
trials from later start years. Our features are the 150 reason-embedding
dimensions, standardized with a scaler fit on the pooled train, validation and
test embeddings. We attempted feature-set variations including additional trial metadata but found that it performed worse than the model trained only on \texttt{why\_stop} embeddings.

\paragraph{Unsupervised step.}
We reduce the standardized embeddings to 30 principal components and fit a
Gaussian mixture with 30 diagonal-covariance components on all 1,561 trials.
Components are initialized from the default $k$-means++ responsibilities of
\texttt{scikit-learn}~\cite{pedregosa2011scikit}, with 3 restarts and the
best-log-likelihood run kept (\texttt{random\_state}=0).
Expectation--maximization~\cite{dempster1977em} stops once the mean
log-likelihood changes by less than $10^{-3}$ between iterations; the selected run
converged after 30 of the permitted 100 iterations.

Each component is then turned into a reason distribution by a
``cluster-then-label'' step. For a given component, we average the label vectors
of the labeled trials, weighting each trial by how strongly the mixture assigns it
to that component. A trial with two reasons contributes half a count to each, so
every trial carries the same total weight regardless of how many reasons it has.
To that average we add a single pseudo-count of the overall reason distribution,
which keeps components holding few or no labeled trials from collapsing onto a
degenerate estimate and instead pulls them toward the global average. A new trial
is scored by taking its component memberships and averaging the corresponding
component distributions.

\paragraph{Supervised base model.}
The supervised half is a multinomial logistic regression over the five reasons
($C = 0.3$, at most 3000 solver iterations). A multi-label trial is handled by
entering it once for each reason it carries, splitting one unit of weight evenly
between those copies, so a two-reason trial contributes half a count to each
rather than counting twice. Classes are then reweighted so that the rare reasons
are not drowned out by Recruitment, which appears on 64\% of trials. We average
this model's predictions with those of a cosine-similarity nearest-neighbour
classifier, giving the regression 70\% of the weight and the neighbours 30\%. The
neighbour model looks at a trial's 10 closest trials and averages their reasons,
weighting each neighbour by the square of its similarity so that close matches
count for much more than distant ones. This average of the two is what we call the
base model.

\paragraph{Self-training.}
We then let the model teach itself~\cite{yarowsky1995unsupervised,triguero2015selflabeled}
over two rounds, using the 1,326 trials with no labels (1,014 from training and
all 312 test trials). In each round every unlabeled trial is scored by mixing the
base model and the mixture, giving the base model 80\% of the say and the mixture
20\%. Any trial whose best reason scores at least 0.8 is treated as if it were
labeled, added to the training pool, and the base model is refit on the enlarged
pool. Starting from 157 labeled trials, 985 of the 1,326 unlabeled trials clear
the bar.

\paragraph{Decision rule.}
Every trial in this corpus stopped for some recorded reason, so a prediction of
``no reason'' is always wrong. We therefore always output the single
highest-scoring reason, and add any further reason scoring at least 0.3.

\paragraph{Model selection.}
The validation set holds only 78 trials, so one trial is worth roughly 1.3 F1
points and small differences between models mean very little. We therefore checked
every choice a second way, by cross-validation over all 235 trials that have real
labels (the 157 labeled training trials plus the 78 validation trials), repeating
the 5-fold split 8 times with different shuffles and averaging. Both numbers are
reported throughout, and where they disagree we follow the cross-validated one.

\subsection{Part 2: predicting trial outcomes}

\paragraph{Baselines.}
Two per-phase baselines, each fit separately on Phase~I, II and III:
a multinomial naive Bayes ($\alpha = 1$) on bag-of-words plus bigram counts of the
eligibility criteria, and a logistic regression ($C = 1$) on criteria TF-IDF
(1--2 grams, $\mathrm{min\_df}=5$, sublinear).

\paragraph{Feature blocks.}
Eight blocks were built over the pooled 10,550 rows: criteria TF-IDF (78,254
features), disease TF-IDF (3,909), drug-name TF-IDF (3,009), ICD codes expanded to
three granularities---chapter letter, 3-character category, 5-character
subcategory (2,382), drug targets (804), drug pathways curated by Reactome~\cite{milacic2024reactome} (1,795), binary drug
actions (31), and five metadata indicators (whether a drug matched, log drug
count, log criteria length, combined-phase flag, early-Phase-1 flag).

\paragraph{Stacked ensemble.}
The improved model combines eight base models: one logistic regression for each
feature block above, plus a naive Bayes and a nearest-neighbour model
($k = 50$) on the criteria text. Each is trained on all three phases at once with
the phase added as an input, so that Phase~I, which has only 1,012 training rows,
can borrow from the much larger Phase~II and III sets. A second logistic
regression then learns how much to trust each of the eight, which is the standard
stacking recipe~\cite{wolpert1992stacked}. So that this combiner is not fooled by
predictions a model has already seen the answer to, each base model's inputs to it
are produced on held-out folds: we split the training trials five ways and predict
each fifth from the other four. The splits keep all rows of a trial together, so
the two rows of a combined-phase trial can never end up on opposite sides of a
split.

\paragraph{Recency weighting.}
Trials have been getting more successful over time---the success rate climbs from
0.38 in 2000 to 0.81 in 2013---and every test trial started later than every
training trial. To stop the model from anchoring on the older, less successful
trials, we make recent trials count for more during training. Each trial's weight
decays exponentially with its age on a three-year time constant, which halves the
weight about every two years, so a trial from the most recent year counts roughly
ten times as much as one from seven years earlier. The weights are applied to both the eight base models and the combiner.
We set the decay rate on two time-ordered holdouts---train on everything up to
2009 and test on 2010--2013, then the harder version training up to 2011 and
testing on 2012--2013---and never on the Kaggle leaderboard.

\paragraph{Error-analysis protocol.}
For the error analysis we take the 452 validation rows that actually failed and
look up each trial's stopping reason from Part~1. 198 of them are found there: 27
with real labels and 171 with labels our Part~1 model predicted. The other 254
never appear in Part~1 because they ran to completion rather than stopping early,
and we keep them as their own ``no recorded reason'' group. Each phase is examined
on its own, so a trial appearing in two phases is never counted twice within one
phase, and a trial with several reasons is counted under each. The error rate we
report is simply the share of these failed trials that the classifier called
successful, meaning it gave them a probability of 0.5 or more. The reason labels
and reason embeddings appear only in this analysis and are never given to the
outcome classifier as inputs.
